% from ega4-8-fr.tex
\oldpage[IV-3]{36}
présentation finie (I, 1.6.2, (i)). Il suffira donc que $T'_j$ soit de
présentation finie sur $U_j$, autrement dit, on peut déjà se borner au cas où
$U_j=S$ et $T_j=T$ est fermé constructible dans $S$.

Soit $S=\operatorname{Spec}(A)$, et considérons $A$ comme limite inductive de
ses sous-$\mathbf Z$-algèbres de type fini, de sorte que
$S=\varprojlim S_\lambda$, où les $S_\lambda$ sont affines et noethériens. En
vertu de (8.3.11), il existe un $\lambda$ et une partie fermée (nécessairement
constructible) $T_\lambda$ de $S_\lambda$ tels que
$T=u_\lambda^{-1}(T_\lambda)$ ($u_\lambda:S\to S_\lambda$ étant le morphisme
canonique). On munira $T_\lambda$ d'une structure de sous-préschéma de
$S_\lambda$ et on prendra $T'=T_\lambda\times_{S_\lambda}S$~; comme
l'immersion canonique $T_\lambda\to S_\lambda$ est de présentation finie
(I, 1.6.1), il en est de même de l'immersion $T'\to S$. Comme $T'$ est affine,
on voit finalement qu'on peut se borner au cas où $T'=S$ est affine.

Alors (8.9.1), avec les mêmes notations, il existe un $\lambda$, un morphisme
de type fini $u_\lambda:X_\lambda\to S_\lambda$ et un
$\mathcal O_{X_\lambda}$-Module cohérent $\mathcal F_\lambda$ tels que $X$
soit isomorphe à $X_\lambda\times_{S_\lambda}S$, $\mathcal F$ à
$\mathcal F_\lambda\otimes_{\mathcal O_{X_\lambda}}\mathcal O_X$. On peut
alors appliquer à $S_\lambda$, $X_\lambda$ et $\mathcal F_\lambda$ le
résultat de (6.9.3) et il y a des sous-préschémas $S'_{\lambda i}$ de
$S_\lambda$, dont les ensembles sous-jacents forment une partition de
$S_\lambda$ et qui sont tels que si l'on pose
$X_{\lambda i}=X_\lambda\times_{S_\lambda}S'_{\lambda i}$,
$\mathcal F_{\lambda i}=\mathcal F_\lambda\otimes_{\mathcal O_{S_\lambda}}
\mathcal O_{S'_{\lambda i}}$ soit plat sur $S'_{\lambda i}$. Les
$S'_i=S'_{\lambda i}\times_{S_\lambda}S$ sont alors des sous-préschémas de
$S$ répondant à la question en vertu de (2.1.4).
\end{proof}

\subsection{Propriétés de permanence des morphismes par passage à la limite projective}
\label{subsection:IV.8.10-fr}

Dans cette section, nous conservons les hypothèses générales et les notations
de (8.8.1).

\begin{proposition}[8.10.1]
\EGAReaderTarget{ega4.c04.prop.8-10-1}{IV.8.10.1-fr}
S'il existe $\lambda$ tel que, pour $\mu\geq\lambda$, $f_\mu$ soit un morphisme
ouvert (resp. universellement ouvert), alors $f$ est un morphisme ouvert
(resp. universellement ouvert).
\end{proposition}

\begin{proof}
Comme $X=X_\lambda\times_{Y_\lambda}Y$, l'assertion relative aux morphismes
universellement ouverts est un cas particulier de (2.4.3, (iii)). Supposons
donc $f_\mu$ ouvert pour $\mu\geq\lambda$~; il suffit de voir que pour tout
ouvert quasi-compact $U$ de $X$, $f(U)$ est ouvert dans $Y$. Or, il existe
$\mu\geq\lambda$ tel que $U=v_\mu^{-1}(U_\mu)$, où $U_\mu$ est un ouvert
quasi-compact dans $X_\mu$ (2.3.11)~; on a alors
$f(v_\mu^{-1}(U_\mu))=w_\mu^{-1}(f_\mu(U_\mu))$ (I, 3.4.8), donc $f(U)$ est
ouvert dans $Y$.
\end{proof}

\begin{corollary}[8.10.2]
\label{IV.8.10.2-fr}
Soit $f:X\to Y$ un morphisme. Pour que $f$ soit universellement ouvert, il
suffit que, pour tout entier $n>0$, si l'on pose
$Y_n=Y\otimes_{\mathbf Z}\mathbf Z[T_1,\ldots,T_n](=\mathbf V_Y^n)$, et
$X_n=X\times_YY_n$, la projection canonique
$f_n=f\times1_{Y_n}:X_n\to Y_n$ soit un morphisme ouvert.
\end{corollary}

\begin{proof}
Pour démontrer que $f$ est universellement ouvert, il suffit de prouver qu'il
en est ainsi de la restriction $f^{-1}(U)\to U$ de $f$ pour tout ouvert affine
$U$ de $Y$~; comme, par hypothèse, si
$U_n=U\otimes_{\mathbf Z}\mathbf Z[T_1,\ldots,T_n]$ est l'image réciproque de
$U$ dans $Y_n$, le morphisme $f_n^{-1}(U_n)\to U_n$, restriction de $f_n$, est
ouvert, on voit qu'on peut se borner au cas où $Y=\operatorname{Spec}(A)$ est
affine. En outre, il suffit aussi évidemment de montrer que pour tout
morphisme $Y'\to Y$, où $Y'=\operatorname{Spec}(A')$ est lui aussi affine,
$f'=f_{(Y')}$ est ouvert. Supposons d'abord que $A'$ soit une $A$-algèbre de
type fini, donc quotient d'une algèbre de polynômes $A[T_1,\ldots,T_n]$~;
alors $Y'$ est un sous-préschéma fermé de $Y_n$ et $f'$ la
